\usepackage{comment}

\begin{comment}
\section{Overview}\label{s:overview}

In this section, we provide a high-level overview of our ABR evaluation framework, focusing on system goals, user-level experience, and interaction between system components. In the latter, the implementation and design choices are explained in Section~\ref{s:design}.

Evaluating ABR algorithms in a fair and realistic manner is a challenge due to diverse network conditions, trace formats, and testbed precision. We aim to address these challenges by creating a reproducible and automated testing environment that allows direct comparison of ABR algorithms as fair as possible, under the same conditions. Puffer's realistic streaming logic inside a real browser and detailed collection of metrics elevate the accuracy of the testing environment.

\parai{System Goals:}

\begin{itemize}
    \item \textbf{Goal 1:} Evaluating the behavior of different ABR algorithms under multiple network conditions.
    \item \textbf{Goal 2:} Analyze application-layer metrics along with transport-layer metrics to observe how each algorithm performs under identical conditions.
    \item \textbf{Goal 3:} Compare ABR performance and rankings using different emulators (Mahimahi and CellReplay).
\end{itemize}

To achieve these goals, it is necessary to build a realistic, controlled, and repeatable setup to reflect the actual playback behavior and ensure a fair comparison.

\parai{User Workflow:}

\begin{itemize}
    \item Select an ABR algorithm supported by Puffer and configure its settings.
    \item Choosing a trace file to replay network conditions (4G or 5G across multiple providers and different signal strengths).
    \item Starting an isolated emulator shell (Mahimahi or CellReplay) that allows the replay of the desired trace.
\end{itemize}

After configuring these settings in the test environment, a single shell command executes a script and launches an automated test run. This includes starting a headless browser session that connects to Puffer's web servers, logs in with specified user and password, navigates to "player" page, and starts streaming under the emulated environment. The script runs a 10-minute playback session, which is as long as the duration of each pre-recorded video. During playback, Puffer logs performance metrics that are used in later analysis.

\parai{System Architecture:}
\begin{itemize}
    \item \textbf{Emulation shell:} Recorded network traces are replayed by Mahimahi or CellReplay, inside an isolated namespace.
    \item \textbf{Browser automation:} A Puppeteer script runs a headless Chrome instance that plays video as a real user would, on a real client.
    \item \textbf{Puffer servers:} Sends video segments from the media server to the web server as chunks with ABR logic while logging application-level metrics.
    \item \textbf{Transport-layer metrics:} Using tcpdump captures, we capture all network packets that are sent, received, or retransmitted during playback. These packets are saved in a PCAP file, to be used in analysis.
    \item \textbf{Analysis tools:} Python scripts process CSV and PCAP files that include metrics from both the application layer and the transport layer, and output them as visualizations, rankings, and summaries.
\end{itemize}

\parai{Trace Collection:}
We use real-world network traces in our experiments.
\begin{itemize}
    \item \textbf{Mahimahi 4G traces:} Mahimahi includes AT\&T, T-Mobile and Verizon LTE profiles in its repository, ready to use.
    \item \textbf{CellReplay 5G traces:} High-resolution recordings from modern 5G networks are included in the CellReplay repository. There are three different T-Mobile and Verizon traces with different signal strengths.
\end{itemize}
To isolate emulator differences, we also run CellReplay's 5G traces inside Mahimahi, to compare if any performance and ranking differ for each ABR algorithm.

\parai{Output and Results:}
\begin{itemize}
    \item Application-layer metrics: SSIM, buffer occupancy, and rebuffering time.
    \item Transport-layer metrics: RTT, TCP retransmissions and delivery rate.
\end{itemize}
These metrics are processed from raw logs and packet traces to generate plots and tables, which are later used for detailed evaluation in Section~\ref{s:evaluation}.
\end{comment}

%%% Local Variables:
%%% mode: latex
%%% TeX-master: "../thesis"
%%% End:
